\documentclass[11pt,a4paper]{ctexart}
\usepackage[a4paper,margin=1.75cm,headheight=14pt]{geometry}
\usepackage{amsmath,amssymb,mathtools}
\usepackage{booktabs,tabularx,longtable,array,multirow}
\usepackage{enumitem}
\usepackage{tikz}
\usetikzlibrary{arrows.meta,positioning,shapes.geometric,fit,calc}
\usepackage[most]{tcolorbox}
\usepackage{xcolor}
\usepackage{fancyhdr}
\usepackage{hyperref}
\usepackage{microtype}
\setlength{\parindent}{2em}
\setlength{\parskip}{0.32em}
\setlist[itemize]{itemsep=0.18em,topsep=0.28em}
\setlist[enumerate]{itemsep=0.18em,topsep=0.28em}
\hypersetup{colorlinks=true,linkcolor=blue!45!black,urlcolor=blue!50!black}

\definecolor{mainblue}{RGB}{45,86,140}
\definecolor{deepblue}{RGB}{30,62,102}
\definecolor{softblue}{RGB}{238,245,252}
\definecolor{softgreen}{RGB}{238,249,243}
\definecolor{softorange}{RGB}{253,246,233}
\definecolor{softred}{RGB}{253,239,239}
\definecolor{softgray}{RGB}{246,247,249}
\definecolor{deepgreen}{RGB}{35,105,75}
\definecolor{deeporange}{RGB}{165,98,22}
\definecolor{deepred}{RGB}{150,55,55}
\definecolor{purple}{RGB}{94,66,140}
\definecolor{softpurple}{RGB}{245,241,251}

\pagestyle{fancy}
\fancyhf{}
\lhead{408 计组 × 操作系统桥梁专题}
\rhead{CPU 与 OS 内核的软硬件协作边界}
\cfoot{\thepage}
\renewcommand{\headrulewidth}{0.4pt}
\setlength{\headheight}{14pt}

\newtcolorbox{corebox}[1][]{colback=softblue,colframe=mainblue,title=#1,fonttitle=\bfseries,boxrule=0.8pt,arc=2mm}
\newtcolorbox{exam}[1][]{colback=softorange,colframe=deeporange,title=#1,fonttitle=\bfseries,boxrule=0.8pt,arc=2mm}
\newtcolorbox{warn}[1][]{colback=softred,colframe=deepred,title=#1,fonttitle=\bfseries,boxrule=0.8pt,arc=2mm}
\newtcolorbox{eng}[1][]{colback=softgreen,colframe=deepgreen,title=#1,fonttitle=\bfseries,boxrule=0.8pt,arc=2mm}
\newtcolorbox{method}[1][]{colback=softgray,colframe=black!45,title=#1,fonttitle=\bfseries,boxrule=0.6pt,arc=2mm}
\newtcolorbox{bridge}[1][]{colback=softpurple,colframe=purple,title=#1,fonttitle=\bfseries,boxrule=0.8pt,arc=2mm}

\newcommand{\kw}[1]{\textbf{#1}}
\newcommand{\code}[1]{\texttt{#1}}

\title{\vspace{-1.2cm}\textbf{CPU 与 OS 内核：软硬件协作边界}\\[0.35em]
\large 系统调用、地址翻译、抢占、DMA 与并发原子性的责任边界}
\author{408 专题手册 · 计组 × 操作系统桥梁篇}
\date{}

\begin{document}
\maketitle

\begin{corebox}[核心模型：硬件机制与内核策略通过接口协作]
\begin{center}
\Large\bfseries CPU/设备提供可被强制执行的局部机制，\par
OS 内核配置这些机制、维护全局状态，并基于目标选择策略。
\end{center}

交接过程：
\[
\boxed{
\text{硬件机制}
\xrightarrow{\text{被 OS 配置}}
\text{受控事件与状态转换}
\xrightarrow{\text{内核处理/策略}}
\text{新的系统状态}
}
\]

分析软硬件边界时依次追问：
\begin{enumerate}
  \item 哪个\kw{物理能力}必须由硬件保证？
  \item OS 预先配置了哪些表、寄存器、权限或队列？
  \item 什么\kw{事件}触发软硬件交接？
  \item 硬件至少保存/检查什么，才能安全把控制权交给内核？
  \item 内核读取哪些状态、执行什么策略、修改哪些数据结构？
  \item 返回前，谁负责恢复可继续执行的体系结构状态？
\end{enumerate}
\end{corebox}

\tableofcontents
\newpage

\section{第一原则：机制、策略与状态归属}

\subsection{把“硬件提供机制，OS 提供策略”展开成三层}
软硬件协作至少包含三层：

\begingroup
\small
\renewcommand{\arraystretch}{1.25}
\setlength{\tabcolsep}{6pt}
\begin{tabularx}{\linewidth}{>{\bfseries}p{0.24\linewidth} X X}
\toprule
层次 & 回答的核心问题 & 典型代表与工程例子\\
\midrule
Hardware Mechanism & 机器\kw{能够且必须可靠地做什么}？ & 特权级检查、MMU 地址翻译、原子 RMW 指令、异常/中断向量表、定时器计数、DMA 搬运\\
Kernel Mechanism & OS 如何把硬件能力封装成可复用系统机制？ & Trap Entry 汇编、页表管理、上下文切换 (\code{switch\_to})、Sleep/Wakeup 队列、驱动框架、互斥锁\\
Policy & 多个合法选择中\kw{应该选哪个}？ & 进程调度算法 (CFS/MLFQ)、页面置换算法 (LRU/CLOCK)、I/O 调度策略、资源回收与降级策略\\
\bottomrule
\end{tabularx}
\endgroup

\begin{bridge}[第三个必须补上的问题：状态到底归谁保存？]
例如一次系统调用横跨三类状态：
\begin{itemize}
  \item CPU 当前寄存器中的\kw{体系结构状态}；
  \item 内核栈或 trap frame 中的\kw{被保存执行状态}；
  \item PCB/TCB、页表、文件表、调度队列中的\kw{内核管理状态}。
\end{itemize}
很多综合题真正考的不是术语，而是“\kw{此刻这份状态在哪里}”。
\end{bridge}

\subsection{六问法：跨计组与 OS 的统一分析模板}
\begin{method}[看到软硬件边界题，固定问六件事]
\begin{enumerate}
  \item \kw{事件是什么？} syscall、page fault、timer interrupt、I/O completion、lock contention？
  \item \kw{硬件先做什么？} 权限检查、最小状态保存、地址转换、原子读改写、发中断？
  \item \kw{OS 预先配置了什么？} 入口地址、页表、定时器、DMA 描述符、等待队列？
  \item \kw{内核接管后做什么？} 判断合法性、记账、调度、分配、唤醒、更新映射？
  \item \kw{什么时候可能发生 context switch？} 进入内核本身不等于切换实体。
  \item \kw{返回时什么状态必须恢复？} PC/SP/权限、页表上下文、寄存器、设备/内存可见性？
\end{enumerate}
\end{method}

\section{控制权与特权级：用户代码怎样高速运行而不失控}

\subsection{先建立架构无关的入口协议}
无论具体 ISA 怎样实现，一次从用户执行进入内核处理，大体都可压成：

\begin{center}

\begingroup
\small
\begin{tikzpicture}[>=Stealth, node distance=1.4cm and 2.4cm]
\tikzset{
  usernode/.style={draw=mainblue, ultra thick, fill=softblue, rounded corners=2.5mm, align=center, minimum width=2.4cm, minimum height=0.95cm, font=\small\bfseries},
  hwnode/.style={draw=purple, ultra thick, fill=softpurple, rounded corners=2.5mm, align=center, minimum width=2.4cm, minimum height=0.95cm, font=\small\bfseries},
  kernnode/.style={draw=deeporange, ultra thick, fill=softorange, rounded corners=2.5mm, align=center, minimum width=2.4cm, minimum height=0.95cm, font=\small\bfseries},
  retnode/.style={draw=deepgreen, ultra thick, fill=softgreen, rounded corners=2.5mm, align=center, minimum width=2.4cm, minimum height=0.95cm, font=\small\bfseries}
}
\node[usernode] (u) {用户执行\\[2pt]\small User Mode};
\node[hwnode, right=of u] (e) {事件识别与硬件入口\\[2pt]\small 查 IDT / 权限 / 切栈};
\node[kernnode, right=of e] (k) {内核入口分发\\[2pt]\small 补全 Context / Handler};

\node[kernnode, below=1.4cm of k] (w) {内核服务/工作\\[2pt]\small Service / Fault / IRQ};
\node[retnode, left=of w] (r) {恢复与返回\\[2pt]\small iret / sysret};

\draw[->, ultra thick, deepblue] (u) -- node[above=2pt, font=\small\bfseries] {syscall / IRQ} (e);
\draw[->, ultra thick, purple] (e) -- node[above=2pt, font=\small\bfseries] {进入 Ring 0} (k);
\draw[->, ultra thick, deeporange] (k) -- node[right=4pt, font=\small\bfseries] {调用处理程序} (w);
\draw[->, ultra thick, deeporange] (w) -- node[below=2pt, font=\small\bfseries] {完成处理} (r);
\draw[->, ultra thick, deepgreen] (r.west) -- node[above=2pt, font=\small\bfseries] {降阶返回 User Mode} ++(-2.4cm,0) |- (u.south);
\end{tikzpicture}
\endgroup
\end{center}

硬件必须至少保证：
\begin{itemize}
  \item 用户不能任意伪造更高特权级；
  \item 进入内核后存在一个可信入口；
  \item 有足够状态使原执行流未来能够继续或被终止；
  \item 非法特权操作能够被捕获，而不是直接破坏全局机器状态。
\end{itemize}

OS 则负责：
\begin{itemize}
  \item 启动阶段配置入口、权限、栈/上下文结构；
  \item 把低级入口统一整理成 trap frame / saved context；
  \item 判断进入原因并调用具体 handler；
  \item 决定返回原执行流、发送信号/终止，还是进入调度器。
\end{itemize}

\subsection{System Call、Exception、Interrupt 的共同点与差异}

\begingroup
\small
\renewcommand{\arraystretch}{1.25}
\setlength{\tabcolsep}{6pt}
\begin{tabularx}{\linewidth}{>{\bfseries}p{0.18\linewidth} >{\raggedright\arraybackslash}p{0.20\linewidth} X X}
\toprule
入口类型 & 与当前指令关系 & 典型触发源头 & 内核接管后首先关心与处理\\
\midrule
系统调用 (Syscall) & 同步、主动发起 & 用户执行 \code{syscall / trap / int} 指令 & 校验调用号、提取参数、检查用户内存合法性、返回结果\\
异常/故障 (Exception) & 同步、被动触发 & 除零、缺页 (\code{\#PF})、非法指令、保护异常 & 异常能否修复并重发指令（如缺页/COW）？还是终止进程？\\
外部中断 (Interrupt) & 异步于当前程序 & Timer 芯片、网卡、磁盘 DMA 控制器 & 哪个硬件设备完成请求？需要把谁从 Blocked 变 Ready？\\
\bottomrule
\end{tabularx}
\endgroup

\begin{exam}[408 核心：进入内核态不等于上下文切换]
系统调用、异常或中断首先造成的是\kw{控制权/特权级转换}。只有内核进一步发现当前执行实体不能继续、应该被抢占，或者主动选择另一个 runnable 实体时，才发生上下文切换。
\end{exam}

\subsection{中断优先级与屏蔽：谁先被响应，和谁能打断谁不是同一问题}

当多个外部中断源竞争 CPU 时，不要把“中断来了”直接画成“CPU 立刻进入某个 ISR”。更完整的接口链是：
\[
\boxed{
\begin{aligned}
\text{Request} &\to \text{Pending} \to \text{Mask/Threshold Eligibility} \to \text{Arbitration} \\
&\to \text{Delivery/Entry} \to \text{Handler} \to \text{EOI/Return}
\end{aligned}
}
\]

这条链同时包含硬件机制与 OS 配置：设备/控制器产生并保存请求状态，控制器根据 mask、优先级或 threshold 判断哪些请求当前有资格递送，再从候选者中仲裁；CPU 完成可信入口后，内核 handler 才开始处理软件语义。

\begingroup
\small
\renewcommand{\arraystretch}{1.25}
\setlength{\tabcolsep}{5pt}
\begin{tabularx}{\linewidth}{>{\bfseries}p{0.22\linewidth} >{\raggedright\arraybackslash}p{0.31\linewidth} X}
\toprule
概念 & 真正回答的问题 & 不要混成什么\\
\midrule
响应优先级 (Response Priority) & 多个\kw{当前可递送}请求同时竞争时，先选择哪一个进入 CPU & 不等于“这个 ISR 执行期间一定能压过所有低优先级 ISR”\\
中断屏蔽 / Priority Threshold & 某个 pending/request 当前是否有资格向 CPU 递送或抢占 & 不等于把硬件的所有优先级编号重新排序；本质上先改变候选集合\\
处理优先级 (Processing Priority) & 当前 handler 执行时，哪些新的中断允许嵌套/抢占它 & 由当前 mask/threshold、硬件优先级机制和 OS 策略共同形成，不应只看中断源的静态名字\\
\bottomrule
\end{tabularx}
\endgroup

在 408 经典中断屏蔽题里，经常把“响应优先级”看作硬件给出的基础仲裁顺序，再由 ISR 设置屏蔽字改变\kw{执行期间允许哪些中断进入}，于是实际处理/完成顺序可以不同于基础响应顺序。现代控制器还可能支持可编程 priority、threshold 与跨核 routing，所以工程上不要把“响应优先级永远固定”写成普遍规律。

\begin{corebox}[控制器的最小职责模型]
中断控制器至少要帮助系统完成四件事：\kw{记录/汇聚请求、筛掉当前不可递送请求、仲裁候选请求、把获胜请求递送/路由到 CPU}。PIC/APIC、不同 ISA 的具体寄存器和应答协议只是这个抽象的不同实现。
\end{corebox}

\begin{exam}[中断优先级/屏蔽题五步]
\begin{enumerate}
  \item 列出当前已经到达并仍然 pending 的中断请求；
  \item 根据题设的 mask/屏蔽字/priority threshold，先确定\kw{当前可递送集合}；
  \item 只在这个集合内部按响应优先级做仲裁，选出下一项；
  \item 进入 ISR 后重新读取该 ISR 的屏蔽规则，判断哪些后续请求允许嵌套抢占；
  \item ISR 返回时回到被打断的执行上下文；若仍有 pending 请求，再按题设规则重新判断，不要凭“原始优先级表”直接写最终完成顺序。
\end{enumerate}
\end{exam}

\begin{warn}[两个常见过度简化]
第一，\kw{masked 不必然等于“请求从此消失”}：请求是否保持 pending、何时重新递送依控制器与触发模式而定，408 题按题设模型判断。第二，\kw{Top/Bottom Half、softirq、workqueue} 是 OS 对中断后续工作的分层机制，不等于“软件把硬件响应优先级改成了另一张固定表”。
\end{warn}

\begin{method}[Bridge 停止规则]
若题目问 IRQ 信号、pending、mask、priority arbitration、vector/routing，主要站在 CPU/中断控制器一侧；一旦 CPU 已进入 handler，继续问“唤醒哪个进程、是否触发调度、设备完成如何改变 Blocked/Ready”，就切回 OS-01/02 与 OS-05。\kw{Bridge 负责交接，不重新定义两侧专题。}
\end{method}

\subsection{x86 例子必须拆成两条路径：不要把 IDT 与 SYSCALL 混成一条}
\begin{warn}[关键误区辨析：传统陷阱门 vs 快速 syscall]
“CPU 查 IDT、自动切内核栈并压入 SS/ESP/EFLAGS/CS/EIP”适合描述\kw{传统 x86 中断/陷阱门跨特权级}的经典路径，但\kw{x86-64 的快速 SYSCALL 指令并不等价于 IDT interrupt gate}。Linux x86-64 也存在多种不同入口，它们的调用约定与栈处理并不完全相同。
\end{warn}

因此手册采用：
\begin{itemize}
  \item \kw{408/架构无关层}：硬件完成可信入口与必要状态转换，内核入口保存其余上下文；
  \item \kw{x86 interrupt/trap gate 例子}：IDT、TSS/IST、硬件形成异常/中断入口 frame；
  \item \kw{x86-64 SYSCALL 例子}：使用专用系统调用入口约定，栈与寄存器整理需要专门入口代码处理，不应套用“IDT 自动压五项”的模板。
\end{itemize}

\begin{eng}[工程边界]
Linux x86 文档明确区分 64-bit system\_call、INT80 compatibility、SYSENTER、普通 interrupt、APIC interrupt 与架构异常等多种入口；不同入口具有不同 calling convention。学习时应先掌握共同机制，再把具体寄存器布局当作架构实现细节。
\end{eng}

\section{虚拟内存：MMU 快速执行翻译，OS 决定映射意味着什么}

\subsection{地址转换是一个“硬件执行、软件定义”的典型契约}
一次普通 load/store 可以抽象为：
\[
VA
\xrightarrow{\text{TLB / page walk}}
PA
\xrightarrow{\text{cache/memory}}
data.
\]

\begingroup
\small
\renewcommand{\arraystretch}{1.25}
\setlength{\tabcolsep}{6pt}
\begin{tabularx}{\linewidth}{>{\bfseries}p{0.22\linewidth} X X}
\toprule
核心关注点 & 硬件 Hardware / MMU 职责 & 软件 OS Kernel 职责\\
\midrule
页表数据结构 & 读取架构规定的页表根指针 (如 CR3) 与 PTE 格式 & 在物理内存中分配页表树，建立、修改与释放映射关系\\
地址翻译过程 & TLB Lookup；未命中时执行硬件 Page Walk & 决定 VPN 到 PFN 的逻辑映射，设置权限/脏位/存在位\\
访问合法性校验 & 执行 PTE 权限检查 (User/Supervisor, R/W)，违规报 Fault & 根据 VMA 区域语义判断是否为 COW、按需分配或非法越界\\
翻译结果加速 & TLB 自动缓存 $VA \rightarrow PA$ 翻译结果 & 当映射更新或释放物理页时，执行 TLB Flush / Shootdown\\
内存页面回收 & ——（硬件不参与策略选择） & 维护全局内存压力，运行 CLOCK / LRU 页面置换与 Swap 写回\\
\bottomrule
\end{tabularx}
\endgroup

\subsection{Page Fault 的本质：硬件发现“当前翻译不能按正常路径完成”}
统一流程：
\[
\text{memory access}
\rightarrow
\text{translation/protection failure}
\rightarrow
\text{fault entry}
\rightarrow
\text{kernel diagnoses cause}
\]
然后可能分叉：
\[
\begin{cases}
\text{合法按需页} & \rightarrow \text{allocate/map}\rightarrow\text{retry}\cr
\text{COW write} & \rightarrow \text{copy/update PTE}\rightarrow\text{retry}\cr
\text{swapped/not resident} & \rightarrow \text{I/O + block}\rightarrow\text{retry}\cr
\text{非法访问} & \rightarrow \text{signal/terminate}
\end{cases}
\]

\begin{corebox}[关键思想]
MMU 不知道“这是不是 COW”“这个页该从哪个文件读取”“是否应该牺牲别的页”。它只知道\kw{当前表项与当前访问是否满足体系结构规则}。\kw{语义解释与策略选择属于 OS。}
\end{corebox}

\subsection{TLB：不要再写“换 CR3 就一定全部清空”}
408 经典模型可以先理解为：切换不同地址空间需要切换页表上下文，并处理旧 TLB translation，带来额外开销。

现代 x86 则存在 PCID 等机制，可以让不同地址空间的 TLB 项带有上下文标识，避免每次切换都简单粗暴地整表清空；具体失效规则取决于硬件功能与内核策略。

\begin{warn}[边界修正]
“TLB hit 一定 1 cycle”“CR3 更新必然完全清空 TLB”“上下文切换会把 CPU Cache 清空”都不应作为普遍结论。408 应掌握\kw{为什么翻译缓存能加速、为什么映射变化必须保持 TLB 一致性}，而不是背某代 CPU 的固定时延。
\end{warn}

\subsection{Accessed/Dirty 位：思想要保留，实现不要绝对化}
在 x86 等体系结构中，硬件可在页表项中维护 accessed/dirty 类状态，OS 可利用这些信息近似判断热度和写回需求。但不同 ISA 对这些位的硬件支持方式并不完全一致。

因此跨架构的正确表达是：
\[
\boxed{\text{硬件/软件共同维护“最近是否访问、是否修改”这类页状态信息}}
\]
它们为页面回收和写回策略提供观测信号，而不是策略本身。

\section{时间虚拟化与抢占：硬件制造调度机会，OS 决定是否换人}

\subsection{OS 如何保证死循环程序也必须交还控制权}
用户程序：
\begin{center}

\begingroup
\large\ttfamily\bfseries
while (1) \{ \}
\endgroup
\end{center}

并发程序并不会主动 syscall。真正保证 OS 最终重新获得 CPU 的关键机制是\kw{可编程定时器 + 外部中断 + 特权入口}。

\begin{center}

\begingroup
\small
\begin{tikzpicture}[>=Stealth, node distance=1.2cm and 1.5cm]
\tikzset{
  usernode/.style={draw=mainblue, ultra thick, fill=softblue, rounded corners=2.5mm, align=center, minimum width=2.5cm, minimum height=0.95cm, font=\small\bfseries},
  hwnode/.style={draw=deeporange, ultra thick, fill=softorange, rounded corners=2.5mm, align=center, minimum width=2.5cm, minimum height=0.95cm, font=\small\bfseries},
  decnode/.style={draw=purple, ultra thick, fill=softpurple, rounded corners=2.5mm, align=center, minimum width=2.5cm, minimum height=0.95cm, font=\small\bfseries},
  retnode/.style={draw=deepgreen, ultra thick, fill=softgreen, rounded corners=2.5mm, align=center, minimum width=2.5cm, minimum height=0.95cm, font=\small\bfseries},
  schnode/.style={draw=deepred, ultra thick, fill=softred, rounded corners=2.5mm, align=center, minimum width=2.5cm, minimum height=0.95cm, font=\small\bfseries}
}
\node[usernode] (u) {进程 A 用户态\\[2pt]\small Running};
\node[hwnode, right=of u] (t) {Timer 倒计时\\[2pt]\small 引爆 IRQ 0};
\node[hwnode, right=of t] (k) {内核 Handler\\[2pt]\small 记账 / slice--};

\node[decnode, below=1.5cm of k] (d) {需要重新调度？\\[2pt]\small need\_resched == 1};
\node[retnode, left=of d] (ret) {No: 返回进程 A\\[2pt]\small 继续执行指令};
\node[schnode, right=of d] (sch) {Yes: 触发 Scheduler\\[2pt]\small Context Switch 换人};

\draw[->, ultra thick, deepblue] (u) -- node[above=2pt, font=\small\bfseries] {硬件打断} (t);
\draw[->, ultra thick, deeporange] (t) -- node[above=2pt, font=\small\bfseries] {进入内核} (k);
\draw[->, ultra thick, deeporange] (k) -- node[right=2pt, font=\small\bfseries] {检查标志} (d);
\draw[->, ultra thick, deepgreen] (d) -- node[above=2pt, font=\small\bfseries] {No} (ret);
\draw[->, ultra thick, deepred] (d) -- node[above=2pt, font=\small\bfseries] {Yes} (sch);
\draw[->, ultra thick, deepgreen] (ret.west) -- node[above=2pt, font=\small\bfseries] {sysret / iret} ++(-1.4cm,0) |- (u.south);
\end{tikzpicture}
\endgroup
\end{center}

硬件提供“\kw{我可以打断正常用户控制流并进入可信内核入口}”；OS 才负责“\kw{这次是否应该换人、换给谁}”。

\subsection{调度不是 timer IRQ 的固定后半段}
408 可以理解成：timer interrupt 提供抢占机会，调度器根据策略选择 runnable task。

更严格地说，具体内核可能先做计时/统计，只设置 reschedule 标记，并在合适的 scheduling point 才真正调用调度器。因此：
\[
\boxed{\text{timer interrupt} \neq \text{必然 context switch}}
\]

\begin{eng}[工程边界：Linux：固定 HZ 不是唯一时间模型]
传统教材喜欢用“每 1ms/10ms 一个时钟中断”建立直觉；现代通用 OS 还广泛使用 local APIC timer、high-resolution timer 与 tickless/nohz 等机制减少无意义周期性 tick。408 核心仍然是：\kw{硬件时间事件创造可抢占点，内核维护调度状态与策略。}
\end{eng}

\subsection{Context Switch 的软硬件边界}
一次真正的切换至少需要：
\begin{enumerate}
  \item 保存旧执行实体恢复所需的 CPU 状态；
  \item 更新旧实体的内核管理状态；
  \item 调度器选择新的 runnable entity；
  \item 必要时切换地址空间上下文；
  \item 恢复新实体的执行状态；
  \item 返回其用户态或继续内核态执行。
\end{enumerate}

不要把某个 Linux 汇编例程保存的“callee-saved registers 列表”背成 CPU 上下文的通用定义：有些寄存器可能已经在 trap frame 中保存，有些由 ABI/入口路径决定。

\section{I/O 与 DMA：硬件搬运、驱动编排、调度隐藏等待}

\subsection{一次块 I/O 的正确分层}
\[
\boxed{
\text{Process}
\rightarrow
\text{Kernel/VFS/Block Layer}
\rightarrow
\text{Driver}
\rightarrow
\text{Device/DMA}
\rightarrow
\text{Interrupt/Completion}
\rightarrow
\text{Wakeup}
}
\]

\begin{center}
\small
\renewcommand{\arraystretch}{1.25}
\begin{tabularx}{\linewidth}{>{\bfseries}p{2.7cm} X X}
\toprule
阶段 & Hardware / Device & Kernel / Driver\\
\midrule
准备 & 暴露 MMIO/PIO 寄存器、DMA capability & 建缓冲/descriptor，做 DMA mapping，写设备寄存器/队列\\
传输 & device/DMA 发起内存总线事务 & 当前线程可阻塞，调度其他 runnable 工作\\
完成 & 更新设备状态，产生 interrupt 或 completion & ISR/driver 确认完成、结束 request、唤醒 waiter\\
继续 & -- & waiter 通常变 Ready/Runnable，未来由 scheduler 再获得 CPU\\
\bottomrule
\end{tabularx}
\end{center}

\subsection{DMA 的真正含义：卸载“逐字搬运”，不是免费传输}
DMA 的核心收益是：CPU 不需要执行一条条 load/store 指令替设备搬完整数据块。

但 DMA 仍然：
\begin{itemize}
  \item 占用 memory/interconnect 带宽；
  \item 可能经过 IOMMU 地址翻译；
  \item 必须处理设备与 CPU 对内存的可见性/一致性；
  \item 在非一致性平台上可能需要 cache synchronization；
  \item 需要 CPU/driver 设置 descriptor、处理完成和错误。
\end{itemize}

\begin{warn}[不要写“DMA 不占 CPU，所以 CPU 一定 100\% 利用率”]
更准确的说法是：\kw{DMA 显著减少 CPU 对数据搬运本身的指令参与}，使 CPU 能执行其他工作；系统最终利用率仍取决于是否存在 runnable workload、内存带宽竞争、驱动开销和设备等待等。
\end{warn}

\subsection{DMA 地址不一定等于 CPU 物理地址}
现代系统可能存在 IOMMU 或其他 DMA 地址转换，因此 driver 交给设备的 DMA address 不一定能被 CPU 直接当作 physical address 使用。

这再次体现：
\[
\boxed{\text{OS 建立映射与权限；硬件执行受控地址访问}}
\]

\section{并发与同步：原子性、内存顺序与“睡眠”必须分层理解}

\subsection{硬件到底为锁提供了什么？}
多核环境中，OS/运行时需要某种\kw{原子 Read-Modify-Write}能力建立竞争点，例如：
\begin{itemize}
  \item compare-and-swap / cmpxchg；
  \item test-and-set / exchange；
  \item fetch-and-add；
  \item LR/SC 一类保留-条件存储机制。
\end{itemize}

硬件还需要给出明确的 memory ordering / fence 语义，使多个核心对共享内存的观察能够按照同步算法要求建立顺序约束。

\begin{bridge}[三个概念不要揉成一个]
\begin{itemize}
  \item \kw{Cache coherence}：多个 cache 对同一 cache line 的副本怎样保持一致；
  \item \kw{Atomicity}：某次 RMW 是否表现为不可被另一核心拆开的单一竞争操作；
  \item \kw{Memory ordering}：不同 load/store 在其他核心看来允许以什么顺序可见。
\end{itemize}
“有 MESI”不等于“程序天然顺序一致”，也不等于“锁算法自动正确”。
\end{bridge}

\subsection{硬件原子指令为什么还不够？}
假设 CAS 只能告诉线程：
\[
\text{lock free?}\rightarrow\text{acquire success/fail}.
\]
如果等待时间很长，失败线程一直 CAS：
\[
\boxed{\text{Spin}}
\]
会浪费 CPU。

因此 OS/运行时进一步提供：
\[
\boxed{
\text{Atomic primitive}
+
\text{wait queue}
+
\text{sleep/wakeup}
+
\text{scheduler}
}
\]
形成 blocking mutex、semaphore、condition variable 等高级同步原语。

\subsection{Fast path / Slow path 再次出现}
\begin{itemize}
  \item 无竞争：用户态原子操作直接成功，快路径；
  \item 有竞争且等待很短：可能 spin；
  \item 有竞争且需要阻塞：进入内核，加入等待队列并 schedule，慢路径。
\end{itemize}

Linux futex 正是理解这种分层的优秀工程例子：常见无竞争路径依赖用户空间原子状态，真正需要等待/唤醒时才通过 futex syscall 进入内核。

\section{五个维度重新统一：同一个“配置 - 触发 - 接管 - 返回”结构}

\begingroup
\small
\renewcommand{\arraystretch}{1.25}
\begin{longtable}{>{\bfseries}p{0.18\linewidth} p{0.28\linewidth} p{0.34\linewidth} p{0.12\linewidth}}
\toprule
核心维度的应用场景 & 硬件 (Hardware) 提供 / 执行的物理机制 & 软件 (OS Kernel) 配置 / 维护 / 决策的逻辑 & 触发界面的关键事件\\
\midrule
\endfirsthead
\toprule
核心维度的应用场景 & 硬件 (Hardware) 提供 / 执行的物理机制 & 软件 (OS Kernel) 配置 / 维护 / 决策的逻辑 & 触发界面的关键事件\\
\midrule
\endhead
特权与系统调用 & 硬件 Privilege Checks；可信入口向量；自动切换内核栈并保护基础断点 & 启动初始化入口向量与权限；在内核栈补全保存 TrapFrame；系统调用分发与返回/调度 & \code{syscall / trap}\\
虚拟内存管理 & MMU 自动硬件 Page Walk；TLB 映射缓存；PTE 权限校验；产生缺页异常 & 分配物理页表树；定义 VMA/对象映射语义；Page Fault Handler；页面回收与 TLB 刷新 & \code{\#PF / TLB miss}\\
分时抢占机制 & 硬件 Timer 倒计时计数；向 CPU 引爆外部硬件中断 (IRQ) & 时间片记账与更新；维护 Ready 队列；触发 \code{need\_resched} 并执行 Scheduler 切换 & \code{timer IRQ}\\
存储与外设 I/O & 暴露 MMIO/PIO 接口；DMA 控制器独立总线数据传输；中断引爆 & 驱动程序装载；配置 DMA 映射 buffer；请求队列编排；线程 Blocked 与唤醒 Wakeup & \code{I/O completion}\\
并发与同步原语 & 硬件原子 RMW 指令 (\code{CMPXCHG/LL-SC})；Memory Fence 屏障；Cache 一致性协议 & 构建 Mutex/Semaphore 高级原语；维护等待队列；阻塞睡眠 \code{sleep()} 与唤醒 \code{wake\_up()} & \code{contention / signal}\\
\bottomrule
\end{longtable}
\endgroup

\begin{corebox}[真正统一的结构]
五个主题表面差异很大，但底层都是：
\[
\boxed{
\begin{aligned}
&\text{OS 先配置规则/状态}\rightarrow\text{硬件高速执行 common case}\rightarrow\text{特殊事件触发边界}\\
&\rightarrow\text{内核接管 慢路径}\rightarrow\text{修复/决策/更新状态}\rightarrow\text{重新回到高速执行}
\end{aligned}
}
\]
这就是理解“CPU × OS 协作”的最强统一视角。
\end{corebox}

\section{综合题：一次阻塞 read() 怎样横跨计组与 OS}

假设用户线程调用：
\begin{verbatim}
read(fd, buf, 4096);
\end{verbatim}
且目标数据不在页缓存，需要真实设备 I/O。

\subsection{第一阶段：控制权交接}
用户代码准备参数并执行系统调用入口。

硬件/体系结构负责：建立受控入口的最低保证。

内核入口负责：补全当前执行流上下文、读取 syscall number/arguments、进入 read 逻辑。

\subsection{第二阶段：资源引用与地址保护}
内核根据：
\[
fd\rightarrow open\ file\ description\rightarrow file/inode
\]
定位打开实例，并验证用户 buffer 是否可以合法写入。

这里已经同时使用了：
\begin{itemize}
  \item 进程资源表；
  \item VFS 抽象；
  \item 虚拟地址空间与保护；
  \item 用户态/内核态边界。
\end{itemize}

\subsection{第三阶段：快路径 失败}
如果 页缓存 hit，数据可能很快返回。

现在假设 miss：
\[
\text{read}
\rightarrow
\text{block layer/driver}
\rightarrow
\text{DMA request}.
\]

当前线程等待设备：
\[
Running\rightarrow Blocked.
\]
内核调度另一个 runnable 实体。

\subsection{第四阶段：设备独立工作}
设备/DMA 在内存与设备之间传输数据；CPU 可以执行其他任务，但 DMA 会使用内存系统与片上/系统互连资源。

完成后设备产生 interrupt/completion。

\subsection{第五阶段：中断、唤醒与重新调度}
CPU 进入 interrupt handler；driver 完成 request bookkeeping，并唤醒等待者：
\[
Blocked\rightarrow Ready.
\]

注意：
\[
\boxed{\text{wakeup} \neq \text{immediately running}}
\]
只有之后 scheduler 选中它：
\[
Ready\rightarrow Running.
\]
线程才真正继续 read 的内核路径，最终返回用户态。

\begin{exam}[一道 read() 已经把两门课串起来]
这条路径同时考察：系统调用入口、CPU 特权级、中断、PCB/线程状态、调度、虚拟内存保护、VFS、I/O、DMA、Cache/页缓存 与阻塞/唤醒。以后遇到跨计组 + OS 综合题，都应该训练成这种“\kw{时间线上谁接管、状态在哪里、谁负责什么}”的推演。
\end{exam}

\section{高频误区：从入门表述回到准确边界}

\begingroup
\small
\renewcommand{\arraystretch}{1.25}
\begin{longtable}{>{\bfseries}p{0.28\linewidth} p{0.64\linewidth}}
\toprule
容易背错的习惯说法 & 更准确稳健的物理与工程模型\\
\midrule
\endfirsthead
\toprule
容易背错的习惯说法 & 更准确稳健的物理与工程模型\\
\midrule
\endhead
“syscall 时 CPU 一定查 IDT、自动压五个寄存器并切内核栈” & 这是把某类 x86 interrupt/trap gate 的经典路径泛化了；快速 syscall 指令拥有不同入口约定。408 掌握“可信入口 + 保存恢复上下文”即可。\\
“进入内核态就是进程切换” & 模式切换只是当前执行流进入内核；只有 scheduler 真正换了执行实体才是 context switch。\\
“换 CR3 就完全清空 TLB” & 经典模型强调切换地址空间会带来翻译缓存开销；现代硬件可用 PCID/ASID 类机制保留不同上下文的 translation。\\
“上下文切换会清空 CPU Cache” & Cache 内容不会因为“进程名变了”自动整体清空；性能损失更多来自 working set 改变、cache/TLB 命中率下降与其他微结构效应。\\
“TLB hit 就是 1 cycle” & 命中很快，但具体 latency 是微架构参数；408 应掌握层次与额外访存次数。\\
“A/D 位都由 MMU 自动维护” & x86 等架构常有硬件维护；跨架构应理解为页访问/脏状态由软硬件共同建立和使用。\\
“CPU 每个固定 1ms/10ms 必然被 IRQ0 打断” & 经典周期 timer 是教学模型；现代系统还有 APIC timer、high-resolution timer 与 tickless 机制。\\
“设置屏蔽字就是把中断响应优先级重新排序” & mask/threshold 首先改变当前可递送的候选集合；候选集合内部仍按题设的仲裁/响应优先级选择。屏蔽规则可以改变实际嵌套与完成顺序，但不能把这两层概念混成一张表。\\
“DMA 不占任何系统资源” & DMA 减少 CPU 搬运指令，但仍消耗 memory/interconnect 带宽，并涉及 mapping、coherency、completion。\\
“I/O 完成中断后等待进程立刻运行” & handler 通常只是使 waiter Ready/Runnable；什么时候 Running 由 scheduler 决定。\\
“有原子指令就有高质量 mutex” & 原子 RMW 只提供竞争原语；阻塞 mutex 还需要等待队列、sleep/wakeup、调度与公平策略。\\
“MESI 保证并发程序顺序正确” & coherence、atomicity、memory ordering 是不同层次；同步算法还需要正确的原子操作与 fence/order 语义。\\
\bottomrule
\end{longtable}
\endgroup

\section{408 跨计组/OS 解题检查表}

\begin{method}[软硬件分工十一问]
\begin{enumerate}
  \item 当前执行在 user 还是 kernel？
  \item 当前 CPU 上是谁？某个进程/线程，还是 interrupt context？
  \item 是同步事件还是异步事件？
  \item 若有多个 IRQ，当前 pending 集合是什么？mask/threshold 后谁仍可递送，响应优先级如何仲裁？
  \item 硬件自动检查/保存/转换了什么？
  \item OS 在事件发生前预先配置了什么？
  \item 内核 handler 要读取哪些寄存器/表项/设备状态？
  \item 当前实体还能继续吗？若不能，是 Ready 还是 Blocked？
  \item 是否真的发生 context switch / page switch / I/O wait？
  \item 这条路径是快路径还是 slow/exception path？
  \item 最终恢复时，哪些体系结构状态和内核不变量必须成立？
\end{enumerate}
\end{method}

\begin{exam}[四类最常见综合题]
\begin{itemize}
  \item \kw{系统调用/中断 + 调度}：何时进入内核？是否必然切换？若多个 IRQ 竞争，谁先进入、谁允许嵌套？
  \item \kw{TLB/Page Table + Page Fault}：硬件翻译到哪一步，OS 从哪一步接管？
  \item \kw{DMA + Interrupt + Process State}：谁搬数据？谁阻塞？谁唤醒？
  \item \kw{Atomic Instruction + Semaphore/Mutex}：硬件保证哪一步原子，OS 为什么还要等待队列和调度？
\end{itemize}
\end{exam}

\section{与相关专题的接口}
\begin{bridge}[责任边界]
\begin{itemize}
  \item 进程与调度专题定义执行实体、状态转换和调度；这里仅说明硬件事件怎样把控制权交给内核。
  \item 虚拟内存专题定义页表、TLB、缺页异常和页面生命周期；这里只保留 MMU 与内核的交接。
  \item 并发专题定义锁、信号量与等待条件；这里只说明原子指令、内存顺序和睡眠机制的分层。
  \item I/O 专题定义请求、DMA、完成与唤醒；这里只说明设备、CPU 与驱动的责任边界。
\end{itemize}
\end{bridge}

\section*{工程边界资料}
\addcontentsline{toc}{section}{工程边界资料}
以下资料用于核对 Linux/x86 的实现边界，不作为 408 必背 API：
\begin{itemize}
  \item Linux x86 Kernel Entries：不同 syscall/interrupt/exception 入口拥有不同 calling convention。\\
  \href{https://www.kernel.org/doc/html/v6.2/x86/entry_64.html}{kernel.org: x86 Kernel Entries}
  \item Linux x86 PTI/PCID：PCID 可降低页表切换时简单全量 TLB flush 的成本。\\
  \href{https://origin.kernel.org/doc/html/latest/arch/x86/pti.html}{kernel.org: x86 PTI}
  \item Linux DMA API：DMA address、coherent/non-coherent mappings 与同步要求。\\
  \href{https://www.kernel.org/doc/html/latest/core-api/dma-api.html}{kernel.org: DMA API}
  \item Linux DMA mapping guide：设备与 CPU 的缓存可见性、DMA buffer 语义。\\
  \href{https://www.kernel.org/doc/html/latest/core-api/dma-api-howto.html}{kernel.org: DMA mapping guide}
  \item Linux futex manual：用户态原子 快路径 与内核 wait/wake 慢路径。\\
  \href{https://man7.org/linux/man-pages/man2/futex.2.html}{man7.org: futex(2)}
\end{itemize}

\begin{corebox}[核心结论]
\begin{center}
\Large\bfseries
CPU/设备负责把规则\textbf{强制执行得足够快、足够可靠}；\par
OS 内核负责决定\textbf{规则怎样配置、状态怎样解释、资源怎样分配、失败以后怎样继续}。
\end{center}
系统能力来自双方稳定而清晰的契约，不是任何一方单独完成全部工作。
\end{corebox}

\end{document}
